\documentclass[11pt,a4paper]{article}
\usepackage[top=1.7cm,bottom=1.5cm,left=1.5cm,right=1.5cm,headsep=8pt,footskip=16pt]{geometry}
\usepackage{enumitem}
\usepackage{amssymb}
\usepackage{array}
\usepackage[table]{xcolor}
\usepackage{fancyhdr}
\usepackage{lastpage}
\usepackage{multicol}
\usepackage{polyglossia}
\setdefaultlanguage[calendar=gregorian, numerals=maghrib]{arabic}
\setotherlanguage{english}
\IfFontExistsTF{Amiri}{%
  \newfontfamily\arabicfont[Script=Arabic, Scale=1.08]{Amiri}
}{%
  \newfontfamily\arabicfont[Script=Arabic, Scale=1.08, Path=../fonts/, Extension=.ttf,
    UprightFont=*-Regular, BoldFont=*-Bold, ItalicFont=*-Italic,
    BoldItalicFont=*-BoldItalic]{Amiri}
}
\IfFontExistsTF{AlBattar}{%
  \newfontfamily\khatt[Script=Arabic, Scale=1.15]{AlBattar}
}{%
  \newfontfamily\khatt[Script=Arabic, Scale=1.15, Path=../fonts/,
    BoldFont=ae_AlBattar.ttf]{ae_AlBattar.ttf}
}
\renewcommand{\headrulewidth}{0.8pt}
\renewcommand{\footrulewidth}{0.8pt}
\fancyhf{}
\rhead{\khatt المدرسة العليا للأساتذة ببوزريعة}
\lhead{\khatt ملحقة بومرداس}
\lfoot{\khatt الأستاذ: د. بن عبد الرزاق يوسف}
\cfoot{\thepage/\pageref{LastPage}}
\pagestyle{fancy}
\parindent 0pt
\setlength{\columnsep}{0.6cm}
\setlength{\columnseprule}{0.4pt}
\setlist[enumerate,1]{label=\textbf{\arabic*.}, leftmargin=*, itemsep=5pt, topsep=2pt}
\begin{document}

\linespread{1.08}\selectfont
\begin{center}
{\LARGE\khatt\textbf{الاختبار السادس في الذكاء الاصطناعي}}\\[4pt]
{\normalsize وضعيات إشكالية ــ أسئلة متعددة الاختيارات ــ الفصول من 1 إلى 6}
\end{center}
\vspace{-2pt}\rule{\textwidth}{0.6pt}\\[2pt]
\textbf{التعليمة:} يتضمن الاختبار خمس وضعيات إشكالية. اقرأ كل وضعية بعناية، ثم اختر التحليل أو التصرف الأنسب
من بين الاختيارات الخمسة؛ لكل وضعية إجابة صحيحة واحدة فقط.
اكتب \textbf{حرف} الإجابة الصحيحة في جدول الإجابات في آخر الاختبار.
لكل وضعية 4 نقاط (المجموع 20 نقطة).
\vspace{4pt}

\noindent
{\large\khatt\textbf{الوضعية 1}}\hfill{\small(الفصل 2)}\par\nopagebreak\vspace{3pt}
\setlength{\fboxsep}{7pt}\colorbox{gray!10}{\parbox{\dimexpr\linewidth-14pt}{طلبت أستاذة من أداة ذكاء اصطناعي عدّ حروف كلمات بيت شعري وتقطيعه عروضيًا، فجاء النص المرافق فصيحًا منظّمًا، لكن عدد الحروف خاطئ في ثلاث كلمات، والتقطيع مختل في تفعيلتين. وتريد شرح ذلك لهم في اجتماع الفريق التربوي. تساءل زملاؤها: كيف تكتب الأداة نصًا سليمًا وتخطئ في عدٍّ بسيط؟ أي تفسير هو الأدق؟}}\par\nopagebreak
\begin{itemize}[leftmargin=1.0cm, labelsep=6pt, itemsep=3pt, topsep=4pt, parsep=0pt]
  \item[\textbf{أ)}] الخطأ يدل على أن الأداة تعمّدت تبسيط الإجابة، لأن النماذج تقدّر مستوى السائل من صياغة سؤاله، فتقدّم له إجابة تقريبية إذا رأت أنه لا يحتاج إلى الدقة في العدّ والتقطيع.
  \item[\textbf{ب)}] النموذج يتقن الصياغة لأنه حفظ دواوين الشعر كاملة، أما الخطأ في العدّ فسببه أن البيت غير موجود في بيانات تدريبه، ولو كان بيتًا مشهورًا لجاء عدّ حروفه وتقطيعه صحيحين بالضرورة، كما يحدث مع الأبيات المتداولة في الكتب المدرسية.
  \item[\textbf{ج)}] النموذج يعدّ الحروف عدًّا دقيقًا كالحاسوب، لكن الحركات وعلامات الضبط تربكه، فإذا حُذفت الحركات من البيت صار العدّ والتقطيع صحيحين دائمًا، دون حاجة إلى أي مراجعة يدوية من الأستاذة.
  \item[\textbf{د)}] النموذج يتنبأ بالمتتالية الأرجح لغويًا، ولذلك يكفي أن تضيف الأستاذة إلى طلبها عبارة «عُدّ بدقة شديدة» ليتحول إلى العدّ الفعلي، فتصبح نتائجه العروضية مضمونة ولا تحتاج إلى أي مراجعة، مهما كان البحر أو طول البيت.
  \item[\textbf{هـ)}] النموذج يعالج النص في وحدات قد تضم عدة حروف، ويتنبأ بالمتتالية الأرجح لغويًا، فيُتقن الصياغة ولا يعدّ الحروف عدًّا فعليًا؛ ولذلك تُراجَع المهام العددية والعروضية يدويًا أو بأداة تنفّذ العدّ برمجيًا.
\end{itemize}
\par\vspace{6pt}
\noindent
{\large\khatt\textbf{الوضعية 2}}\hfill{\small(الفصل 3)}\par\nopagebreak\vspace{3pt}
\setlength{\fboxsep}{7pt}\colorbox{gray!10}{\parbox{\dimexpr\linewidth-14pt}{طلب أستاذ في طلب واحد إعداد وحدة تعليمية كاملة: نص قرائي، وأسئلة فهم، ودرس نحو مستخرج من شواهد النص، وتمارين، ووضعية إدماجية، وشبكة تقويم. جاءت النتيجة طويلة لكنها سطحية، والدرس النحوي لا يرتبط بالنص. ويريد وحدة مترابطة قابلة للتنفيذ. وقد سبق له أن جرّب ذلك مع وحدة أخرى بالنتيجة نفسها. ما الاستراتيجية الأنجع؟}}\par\nopagebreak
\begin{itemize}[leftmargin=1.0cm, labelsep=6pt, itemsep=3pt, topsep=4pt, parsep=0pt]
  \item[\textbf{أ)}] يقسّم المهمة إلى طلبات متتابعة، لكنه يطلب الدرس النحوي والتمارين أولًا، ثم يطلب نصًا تُحشر فيه شواهد الدرس، لأن النص أسهل العناصر، ويمكن تكييفه مع ما سبقه في كل الأحوال.
  \item[\textbf{ب)}] يعيد الطلب الواحد نفسه مع إضافة «كن عميقًا ومترابطًا جدًا»، لأن الصفات التي تصف الجودة المطلوبة تكفي لرفع عمق الإجابة وربط الدرس النحوي بالنص، دون تغيير بنية الطلب أو تقسيمه، فالمشكلة في نبرة الطلب لا في بنيته.
  \item[\textbf{ج)}] يقسّم المهمة إلى طلبات متتابعة في المحادثة نفسها: النص أولًا ثم مراجعته، ثم الأسئلة انطلاقًا منه، ثم الدرس النحوي من شواهده، وهكذا، لأن كل خطوة تُبنى على مُخرَج سابق راجعه وثبّته.
  \item[\textbf{د)}] يقسّم المهمة إلى ستة طلبات يرسلها في ست محادثات منفصلة ومتزامنة، لأن فصل المحادثات يمنع تداخل المهام، ثم يجمع المخرجات الستة في ملف واحد، فتأتي الوحدة مترابطة تلقائيًا.
  \item[\textbf{هـ)}] يطلب الوحدة كاملة من أداة أخرى أحدث، لأن السطحية دليل على ضعف الأداة لا على طريقة الطلب، والأدوات الأحدث تنجز المهام المركبة في طلب واحد بعمق وترابط مضمونين دون مراجعة، كما يؤكد كثير من المستخدمين.
\end{itemize}
\par\vspace{6pt}
\noindent
{\large\khatt\textbf{الوضعية 3}}\hfill{\small(الفصل 4)}\par\nopagebreak\vspace{3pt}
\setlength{\fboxsep}{7pt}\colorbox{gray!10}{\parbox{\dimexpr\linewidth-14pt}{اقترح نادي المؤسسة إنتاج تسجيل بصوت شاعر جزائري راحل يلقي قصيدة له لم يُسجَّل صوته وهو يلقيها، باستعمال أداة استنساخ صوتي تتدرّب على تسجيلات قديمة له، ثم نشره على صفحة المؤسسة بمناسبة وطنية. ويرى أعضاء النادي أن الهدف نبيل. وطلب رئيس النادي رأي أستاذ اللغة العربية. ما الموقف الأسلم أخلاقيًا؟}}\par\nopagebreak
\begin{itemize}[leftmargin=1.0cm, labelsep=6pt, itemsep=3pt, topsep=4pt, parsep=0pt]
  \item[\textbf{أ)}] لا حرج في إنتاجه ونشره دون تنبيه، لأن القصيدة للشاعر نفسه، فالصوت المولَّد لا يضيف إليه قولًا لم يقله، ولا يحتاج النادي إلى إذن من أحد ما دام الهدف تربويًا ووطنيًا نبيلًا.
  \item[\textbf{ب)}] لا يُنتَج التسجيل إلا بإذن أصحاب الحقوق، ويُصرَّح بوضوح بأنه صوت مولَّد لا تسجيل حقيقي، مع الالتزام بنص القصيدة حرفيًا؛ والأسلم تربويًا أن يلقيها تلميذ بصوته.
  \item[\textbf{ج)}] يُنتَج التسجيل ويُنشر مع التصريح بأنه مولَّد، دون حاجة إلى إذن أصحاب الحقوق، لأن التصريح وحده يرفع كل إشكال أخلاقي وقانوني، مهما كان استعمال صوت الشخص ومهما اتسع انتشار المنشور.
  \item[\textbf{د)}] يُنتَج التسجيل بإذن أصحاب الحقوق، ويُنشر دون تنبيه إلى أنه مولَّد، لأن الإذن يجعل الصوت المستنسخ في حكم التسجيل الحقيقي، والتنبيه قد يُضعف الأثر الوجداني للمناسبة عند الجمهور.
  \item[\textbf{هـ)}] يُرفض المشروع، لأن كل استعمال لأدوات الصوت المولَّد ممنوع في المؤسسات التربوية، حتى قراءة نصوص الكتاب المدرسي بأصوات اصطناعية لا تحاكي شخصًا معيّنًا، ولو صُرّح بطبيعتها بوضوح.
\end{itemize}
\par\vspace{6pt}
\noindent
{\large\khatt\textbf{الوضعية 4}}\hfill{\small(الفصل 5)}\par\nopagebreak\vspace{3pt}
\setlength{\fboxsep}{7pt}\colorbox{gray!10}{\parbox{\dimexpr\linewidth-14pt}{نقل طالب فقرات من مقال محكَّم، ثم طلب من أداة ذكاء اصطناعي «إعادة صياغتها حتى لا يكشفها برنامج كشف الانتحال»، فجاءت الصياغة مختلفة الألفاظ تمامًا، ولم يرصدها البرنامج. ثم أدرجها في مذكرته دون إحالة إلى المقال. واكتشف المشرف الأمر عند قراءته المقال نفسه. كيف يُقيَّم عمله في ضوء أخلاقيات البحث؟}}\par\nopagebreak
\begin{itemize}[leftmargin=1.0cm, labelsep=6pt, itemsep=3pt, topsep=4pt, parsep=0pt]
  \item[\textbf{أ)}] هو انتحال رغم تغيير الألفاظ، لأن الأفكار والبناء الحجاجي لصاحب المقال، والإحالة واجبة على الأفكار لا على الألفاظ وحدها؛ والصواب أن يوثّق المقال، ويقتبس نصًا أو يلخّص بأسلوبه مع الإحالة.
  \item[\textbf{ب)}] هو انتحال، والحل أن يطلب من الأداة صياغة أبعد عن الأصل، لأن الانتحال يزول عندما تبلغ نسبة الاختلاف في الألفاظ حدًا معينًا، ولا تحتاج الأفكار المنقولة إلى توثيق بعد ذلك.
  \item[\textbf{ج)}] عمله مقبول بشرط أن يذكر في مقدمة المذكرة أنه استعمل أداة ذكاء اصطناعي في إعادة الصياغة، لأن الإفصاح عن الأداة يغني عن الإحالة إلى المقال الأصلي الذي أُخذت منه الأفكار، كما تنص على ذلك بعض أدلة البحث.
  \item[\textbf{د)}] الخلل الوحيد في استعمال الأداة، فلو أعاد الصياغة بنفسه لصار عمله سليمًا حتى دون إحالة، لأن إعادة الصياغة الشخصية تكفي لتحويل أفكار الآخرين إلى إنتاج أصيل يملكه الطالب.
  \item[\textbf{هـ)}] عمله سليم، لأن الانتحال يتعلق بنقل الألفاظ حرفيًا، وما دام البرنامج لم يرصد أي تطابق، فالنص الجديد صار ملكًا للطالب، ولا تلزمه الإحالة إلى المقال الذي استوحى منه أفكاره.
\end{itemize}
\par\vspace{6pt}
\noindent
{\large\khatt\textbf{الوضعية 5}}\hfill{\small(الفصل 6)}\par\nopagebreak\vspace{3pt}
\setlength{\fboxsep}{7pt}\colorbox{gray!10}{\parbox{\dimexpr\linewidth-14pt}{طلبت أستاذة من أداة توليد صور ملصقًا عن «أقسام الكلمة» يتضمن رسومًا توضيحية وعناوين عربية. جاءت الرسوم جميلة، لكن الكلمات العربية داخل الصورة مشوّهة: حروف منفصلة، وكلمات غير مفهومة، وبعضها معكوس الاتجاه. وتريد تعليق الملصق في القسم. وقد جرّبت إعادة الطلب مرتين دون تحسّن يُذكر. ما التصرف الأنسب؟}}\par\nopagebreak
\begin{itemize}[leftmargin=1.0cm, labelsep=6pt, itemsep=3pt, topsep=4pt, parsep=0pt]
  \item[\textbf{أ)}] تعتمد الملصق كما هو، وتطلب من التلاميذ تصحيح الكلمات المشوهة نشاطًا إملائيًا، لأن الملصق المعلّق في القسم يُستحسن أن يتضمن أخطاء مقصودة تنمّي يقظة التلاميذ اللغوية طوال السنة.
  \item[\textbf{ب)}] تعيد الطلب مرات كثيرة حتى تظهر الكلمات سليمة، ثم تعتمد أول صورة مقبولة دون تدقيق حرفي، لأن الكلمات التي تبدو سليمة للوهلة الأولى تكون صحيحة الإملاء بالضرورة في الصور المولَّدة.
  \item[\textbf{ج)}] تكتب العناوين في الطلب بالحروف اللاتينية، لأن الأداة ستحوّلها تلقائيًا إلى حروف عربية سليمة متصلة داخل الصورة، فالمشكلة في لغة الطلب لا في قدرة الأداة على رسم الحروف العربية.
  \item[\textbf{د)}] تستعمل الأداة لتوليد الرسوم وحدها دون نصوص، ثم تضيف العناوين العربية بنفسها في برنامج تصميم أو معالجة نصوص، لأن أدوات توليد الصور كثيرًا ما تُخفق في رسم الحروف العربية المتصلة.
  \item[\textbf{هـ)}] تتخلى عن الصور المولَّدة نهائيًا، لأن تشوّه الكلمات دليل على أن الرسوم نفسها غير دقيقة علميًا، فلا يمكن الاستفادة من أي جزء من الصورة حتى لو أُضيفت النصوص يدويًا بعد ذلك.
\end{itemize}
\par\vspace{6pt}
\begin{center}
{\khatt\textbf{جدول الإجابات}}\\[3pt]
\renewcommand{\arraystretch}{1.6}
\begin{tabular}{|>{\centering\arraybackslash}m{1.9cm}|*{5}{>{\centering\arraybackslash}m{1.2cm}|}}
\hline
\textbf{رقم السؤال} & \textbf{1} & \textbf{2} & \textbf{3} & \textbf{4} & \textbf{5} \\ \hline
\textbf{الإجابة} &  &  &  &  &  \\[8pt] \hline
\end{tabular}

\vspace{8pt}
\renewcommand{\arraystretch}{2.2}
\begin{tabular}{|>{\centering\arraybackslash}m{7.5cm}|>{\centering\arraybackslash}m{3.2cm}|>{\centering\arraybackslash}m{3.2cm}|}
\hline
\textbf{الاسم واللقب} & \textbf{الفوج} & \textbf{العلامة} \\ \hline
 & & \hfill /20 \\ \hline
\end{tabular}
\end{center}
\end{document}
